\section{Diagnostic Probes and Policy Selection}

We freeze the MLLM backbone and train LoRA parameters $\Delta$ \citep{hu2021lora}. For a chart image $x$, question $q$, and answer $y$, the adapted model defines $p_{\theta,\Delta}(y\mid x,q)$ with frozen $\theta$. The methods below form a probe ladder rather than a single proposed fix. Directional adapters test whether scale preference can be moved; compute-matched controls test whether the movement is due to the scale objective; clean-direct controls test whether it is safe without answer tuning; blind sampling tests whether the correct answer is generated at all; and selectors test whether candidate evidence can be turned into a deployable regime decision.

\begin{figure}[t]
\centering
\small
\setlength{\fboxsep}{4pt}
\fbox{\begin{minipage}{0.95\linewidth}
\textbf{Chart + question}
$\rightarrow$
\textbf{Base answer} (e.g., \textit{0.61})\\[-0.15em]
\rule{\linewidth}{0.3pt}\\[-0.2em]
\textbf{Candidate pool:}
base, self-check answer, directional adapter answers, blind samples, explicit scale variants.\\[-0.15em]
\textbf{Regime evidence:}
question form, answer magnitude, percent/unit tokens, axis labels, tick values, candidate likelihoods.\\[-0.15em]
\textbf{Policy decision:}
base, self-check, prefer-large, prefer-small, decimal guard, or unit guard.\\[-0.15em]
\rule{\linewidth}{0.3pt}\\[-0.1em]
\begin{tabular}{@{}l@{}}
\emph{Prediction too small:} $0.61 \rightarrow 61$ \\
\emph{Prediction too large:} $61 \rightarrow 0.61$ \\
\emph{True small decimal:} keep $0.0593$ \\
\emph{Unit mismatch:} millions vs. raw values \\
\end{tabular}
\end{minipage}}
\caption{Regime-aware diagnostic view. The experiments test whether scale-related candidates can be generated and, more importantly, whether a deployable policy can select the scale regime required by the chart and question.}
\label{fig:regime_pipeline}
\end{figure}

\paragraph{Directional adapters as probes.}
For numeric training examples we construct scale-related pairs $(y^+,y^-)$ in which $y^-$ is $y^+$ rescaled by a factor of $100$ or $1/100$ (with percent signs handled as in \S\ref{sec:strict}; full pair-construction rules, rounding, and sign handling are in Appendix~\ref{app:pairs}). \emph{Prefer-large} pairs take the larger-scale value as $y^+$; \emph{prefer-small} pairs take the smaller. Tables retain the training-time labels \emph{large-positive}/\emph{small-positive}, which name the same preferences. The names describe the training preference, not the observed error direction: an answer that is too small requires a prefer-large correction, and vice versa. We also evaluate balanced mixtures, all mined pairs, and real-conflict pairs mined from model disagreements.

\paragraph{Objective.}
We compute mean token negative log-likelihood $\mathrm{NLL}(y\mid x,q)$ and minimize
\begin{equation}
\begin{aligned}
\mathcal{L}_{\mathrm{scale}}
=\mathrm{softplus}\big(&-\beta[
\mathrm{NLL}(y^-) \\
&-\mathrm{NLL}(y^+)-m]\big),
\end{aligned}
\end{equation}
where $m$ is a margin and $\beta$ controls sharpness. The loss encourages the preferred answer to have lower NLL than its scale-related alternative.

\paragraph{Adaptation regimes and controls.}
The main positive result uses an \emph{answer-tuned} starting point: scale adaptation continues from an answer-only LoRA adapter in the Qwen2.5-VL-7B pipeline for 300 additional steps. Because any continuation could help simply through more optimization, we add \emph{compute-matched controls} that continue the same answer-tuned adapter for the same 300 steps with (a)~ordinary answer examples, (b)~random or shuffled numeric negatives, and (c)~non-directional balanced pairs with matched answer frequency. We also test \emph{clean direct} scale-only tuning from the raw backbone as a safety control. At inference, main results use the same deterministic short-answer decoding template as the baseline, without reranking or calculators.

\paragraph{Blind-generation oracle.}
To test whether candidate oracles are high only because the pool contains deterministic scale transforms of the base answer, we also draw $20$ stochastic samples from the answer-tuned model with no scale prompt and no transformation, and report single-sample accuracy against pass@20. This separates ``the model can generate the right scale'' from ``our transformation rules reproduce the mined relation.''

\paragraph{Label-free policy selection.}
\label{sec:router_method}
ScaleQA-Natural suggests that no single direction is safe across all regimes. We therefore evaluate selectors that choose among base answering, self-check prompting, unit/axis guards, and directional adapters, using question text, base prediction text, numeric features, and chart-context strings extracted from the image (title, axis labels, tick values, legends, units). We call a selector \emph{label-free} when it uses no gold subtype or direction field; this is a weaker guarantee than an untouched test set, and we report selector results under a fixed partition and leave-source-out evaluation accordingly (\S\ref{sec:selector}). Oracle subtype, policy-pool, and candidate-pool results are upper bounds whose role is to show how much headroom remains if regime detection were solved.
